\documentclass{article}

\usepackage[preprint]{neurips_2025}

\usepackage[utf8]{inputenc}
\usepackage[T1]{fontenc}
\usepackage{courier}
\usepackage{hyperref}
\usepackage{url}
\usepackage{booktabs}
\usepackage{amsfonts}
\usepackage{amsmath}
\usepackage{graphicx}
\usepackage{microtype}
\usepackage{xcolor}
\usepackage{lineno}

\linenumbers

\title{Reproducible Training-Free Inference Acceleration on Pythia-70M}

\author{%
  Ruoyu Fu \\
  NLP Course Final Project \\
  Code: \href{https://github.com/lilfry09/Training-free-LM-inference-acceleration}{GitHub repository}
}

\begin{document}
\hypersetup{hidelinks}

\maketitle

\begin{abstract}
Efficient inference is a practical bottleneck for autoregressive language models because attention computation and the KV cache grow with context length. This work studies training-free inference modifications on Pythia-70M: a reduced-KV grouped-query attention (GQA) conversion and two KV cache compression strategies implemented with KVPress, Knorm and StreamingLLM. We evaluate exact teacher-forced perplexity, cached next-token perplexity, time to first token (TTFT), time per output token (TPOT), and throughput on WikiText-103 and a PG-19 sample. KVPress improves throughput in several CPU settings, especially WikiText at 128/512-token contexts, while the training-free GQA conversion stores only 2 KV heads but severely degrades perplexity and does not provide a usable speed-quality tradeoff. The experiments can be reproduced with the scripts included in the repository.
\end{abstract}

\section{Introduction}

Autoregressive decoding repeatedly evaluates a transformer while maintaining a key-value (KV) cache for previous tokens. This cache avoids recomputing attention states, but its memory and attention cost scale with sequence length. Recent efficient-inference work therefore attacks the problem from several directions: using fewer KV heads \citep{ainslie2023gqa}, more IO-aware attention kernels \citep{dao2022flashattention}, or compressing the KV cache during generation \citep{xiao2023streamingllm,kvpress2024}.

This project focuses on methods that can be reproduced without training a new model. We use Pythia-70M \citep{biderman2023pythia}, a small GPT-NeoX model, as the common backbone. The assignment asks for both quality and speed evaluation on datasets such as WikiText and PG-19, so we report two complementary quality metrics. First, exact perplexity is computed in the standard teacher-forced way and is not affected by KVPress, because KVPress acts on generation caches rather than full-context forward passes. Second, cached perplexity evaluates next-token prediction after a prefill pass whose KV cache may be compressed. This second metric is the appropriate quality signal for KVPress in this implementation.

The main contribution is a reproducible comparison under a CPU-only course environment. The results show that reduced-KV GQA is not a reliable drop-in conversion for a model trained with standard multi-head attention, while KV cache compression can help on some contexts but is sensitive to runtime overhead.

\section{Methods}

\paragraph{Baseline.}
The baseline loads the local \texttt{EleutherAI/pythia-70m} checkpoint through Hugging Face Transformers and uses the default scaled dot-product attention (SDPA) implementation. No weights are changed.

\paragraph{Reduced-KV GQA conversion.}
Pythia-70M uses 8 attention heads. In the GQA variant, we compute the original query, key, and value states, average the key/value heads into 2 grouped KV heads, and run custom grouped attention without materializing repeated K/V states. The dynamic cache therefore stores tensors with shape \texttt{[batch, 2, seq, head\_dim]} rather than 8 KV heads. This is still a training-free conversion rather than the fine-tuned GQA recipe of \citet{ainslie2023gqa}, so quality degradation is expected.

\paragraph{KVPress cache compression.}
We evaluate two KVPress modes at compression ratio 0.5. Knorm keeps tokens with larger key norm scores. StreamingLLM keeps attention-sink tokens and recent tokens, following the observation that early sink tokens are important for stable long-context decoding \citep{xiao2023streamingllm}. In both cases, compression is applied after the prefill stage via forward hooks on GPT-NeoX attention layers. The implementation records per-layer cache lengths before and after compression, e.g., 512 cached positions are compressed to 256.

\paragraph{Metrics.}
Exact PPL is computed over the first 4096 tokens using 512-token chunks. Cached PPL first prefills 512 context tokens, optionally compresses the resulting cache, and then scores the next 256 ground-truth tokens autoregressively with the cache. Speed is measured with manual greedy decoding for 64 new tokens at context lengths 128, 512, and 1024. Each setting uses one warm-up and three measured runs. TTFT is wall time from prefill start to selecting the first token; TPOT is decode time divided by the remaining generated tokens; throughput is generated tokens per end-to-end second.

\section{Experimental setup}

We evaluate the first 4096 tokens from the test split of WikiText-103 \citep{merity2016wikitext} and from one PG-19 test sample \citep{rae2020compressive}, matching the course requirement that PG-19 can be evaluated with a single long sample. The local environment is CPU-only with PyTorch 2.6.0, Transformers 4.56.2, Datasets 3.6.0, and KVPress 0.5.3. Since CUDA is unavailable, FlashAttention is not used as a main experimental condition.

The exact reproduction command is:
\begin{verbatim}
cd finalproj
.\scripts\run_matrix.ps1
\end{verbatim}
This script runs baseline, GQA, KVPress-Knorm, and KVPress-StreamingLLM on both datasets and regenerates \texttt{comparison\_quality.*} and \texttt{comparison\_speed.*}.

\section{Results and discussion}

Table~\ref{tab:quality} summarizes quality. The exact-PPL column should be used to compare baseline and GQA. The cached-PPL column should be used to compare KVPress cache-compression behavior, because it scores continuation tokens after compressed prefill. GQA causes a large degradation: WikiText exact PPL rises from 63.72 to 2289.96, and PG-19 exact PPL rises from 33.57 to 797.49. This confirms that converting a pretrained MHA checkpoint into 2 cached KV heads by simple averaging is not quality-preserving.

\begin{table}[t]
  \caption{Quality metrics. Exact PPL is standard teacher-forced PPL; cached PPL scores 256 continuation tokens after a 512-token prefill cache.}
  \label{tab:quality}
  \centering
  \small
  \setlength{\tabcolsep}{4pt}
  \begin{tabular}{llrrrr}
    \toprule
    Dataset & Method & Exact PPL & Cached PPL & KV events & Avg. comp. \\
    \midrule
    WikiText & Baseline & 63.72 & 13.09 & 0 & -- \\
    WikiText & GQA-2KV & 2289.96 & 3046.28 & 0 & -- \\
    WikiText & KVPress-Knorm & 63.72 & 52.44 & 78 & 0.50 \\
    WikiText & KVPress-StreamingLLM & 63.72 & 61.30 & 78 & 0.50 \\
    PG-19 & Baseline & 33.57 & 29.60 & 0 & -- \\
    PG-19 & GQA-2KV & 797.49 & 781.27 & 0 & -- \\
    PG-19 & KVPress-Knorm & 33.57 & 31.76 & 78 & 0.50 \\
    PG-19 & KVPress-StreamingLLM & 33.57 & 29.49 & 78 & 0.50 \\
    \bottomrule
  \end{tabular}
\end{table}

The WikiText baseline has much lower cached PPL than full-window exact PPL. To check whether this was a cache-scoring artifact, we recomputed ordinary teacher-forced PPL on the identical continuation slice: the first 512 tokens were used as context and the next 256 tokens were scored. The same-slice teacher-forced PPL was 13.088, while cached autoregressive PPL was 13.089. This small difference indicates that cached scoring matches standard teacher forcing on the same tokens. The gap from the 63.72 exact PPL is therefore caused mainly by the evaluation window: exact PPL averages over 4095 tokens, whereas cached PPL uses a fixed local continuation. We therefore compare cached PPL only across methods on the same cached-PPL slice, not against exact PPL as an absolute model-quality number.

Table~\ref{tab:speed} reports mean throughput changes relative to the baseline over three measured runs. On WikiText, KVPress improves 128/512-token generation, but both KVPress variants slow down at 1024 tokens. StreamingLLM gives +35.3\% and +50.3\% at 128 and 512 tokens, respectively. On PG-19, StreamingLLM is close to baseline at 512 tokens and gives a modest +8.7\% at 1024 tokens. GQA stores fewer KV heads but is not a useful accelerator because its quality collapses and its speed is inconsistent.

\begin{table}[t]
  \caption{Throughput change relative to baseline. Positive values indicate faster generation.}
  \label{tab:speed}
  \centering
  \small
  \setlength{\tabcolsep}{4pt}
  \begin{tabular}{llrrr}
    \toprule
    Dataset & Method & 128 ctx & 512 ctx & 1024 ctx \\
    \midrule
    WikiText & GQA-2KV & -30.7\% & +23.1\% & -42.7\% \\
    WikiText & KVPress-Knorm & +31.3\% & +49.7\% & -14.3\% \\
    WikiText & KVPress-StreamingLLM & +35.3\% & +50.3\% & -17.5\% \\
    PG-19 & GQA-2KV & -6.1\% & -8.8\% & -12.3\% \\
    PG-19 & KVPress-Knorm & -12.4\% & -14.7\% & -14.0\% \\
    PG-19 & KVPress-StreamingLLM & -10.6\% & -3.0\% & +8.7\% \\
    \bottomrule
  \end{tabular}
\end{table}

The raw TTFT, TPOT, end-to-end time, throughput means, standard deviations, and per-run values are included in \texttt{outputs/comparison\_speed.*}. The CPU results are mixed: compression helps in some settings but not all, and small deltas should not be over-interpreted. A GPU implementation with fused kernels and larger models would be a more favorable setting for cache compression and reduced-KV attention.

\section{Conclusion}

The most reliable positive results are KVPress on WikiText at 128/512-token contexts and StreamingLLM on PG-19 at 1024 tokens, where the 3-run mean throughput improves while the cache is explicitly compressed by 50\%. Reduced-KV GQA is not successful despite storing only 2 KV heads: it sharply worsens perplexity and does not consistently accelerate generation. Overall, training-free KV cache compression is a practical direction for inference acceleration, but its benefit depends on implementation overhead, dataset/runtime behavior, and model size.

{\small

\begin{thebibliography}{9}

\bibitem[Biderman et~al.(2023)]{biderman2023pythia}
Biderman, S., Schoelkopf, H., Anthony, Q., Bradley, H., O'Brien, K., Hallahan, E., Khan, M.~A., Purohit, S., Prashanth, U.~S., Raff, E., Skowron, A., Sutawika, L., and Van~Der~Wal, O.
\newblock Pythia: A suite for analyzing large language models across training and scaling.
\newblock \emph{International Conference on Machine Learning}, 2023.

\bibitem[Ainslie et~al.(2023)]{ainslie2023gqa}
Ainslie, J., Lee-Thorp, J., de~Jong, M., Zemlyanskiy, Y., Lebron, F., and Sanghai, S.
\newblock GQA: Training generalized multi-query transformer models from multi-head checkpoints.
\newblock \emph{Empirical Methods in Natural Language Processing}, 2023.

\bibitem[Dao et~al.(2022)]{dao2022flashattention}
Dao, T., Fu, D.~Y., Ermon, S., Rudra, A., and Re, C.
\newblock FlashAttention: Fast and memory-efficient exact attention with IO-awareness.
\newblock \emph{Advances in Neural Information Processing Systems}, 2022.

\bibitem[Xiao et~al.(2023)]{xiao2023streamingllm}
Xiao, G., Tian, Y., Chen, B., Han, S., and Lewis, M.
\newblock Efficient streaming language models with attention sinks.
\newblock arXiv:2309.17453, 2023.

\bibitem[NVIDIA(2024)]{kvpress2024}
NVIDIA.
\newblock KVPress: KV cache compression for efficient LLM inference.
\newblock \url{https://github.com/NVIDIA/kvpress}, 2024.

\bibitem[Merity et~al.(2016)]{merity2016wikitext}
Merity, S., Xiong, C., Bradbury, J., and Socher, R.
\newblock Pointer sentinel mixture models.
\newblock arXiv:1609.07843, 2016.

\bibitem[Rae et~al.(2020)]{rae2020compressive}
Rae, J.~W., Potapenko, A., Jayakumar, S.~M., and Lillicrap, T.~P.
\newblock Compressive transformers for long-range sequence modelling.
\newblock \emph{International Conference on Learning Representations}, 2020.

\end{thebibliography}
}

\end{document}
